\title{\vspace{-1.4cm}\textbf{Valutazione di OpenTelemetry\\per il monitoraggio di workload mixed-criticality}\\[.3em]
\large Campagna sperimentale su Linux PREEMPT-RT}
\author{Benito Cervone \and \textit{(secondo componente del gruppo)}}
\date{\today}

\begin{document}
\maketitle
\thispagestyle{empty}

\begin{center}
\begin{minipage}{0.9\textwidth}
\small\itshape
Questo documento riporta il lavoro svolto per valutare se e in che misura OpenTelemetry sia in
grado di monitorare workload real-time a criticit\`a mista, con particolare attenzione a due
domande: se la pipeline di telemetria sappia \emph{prioritizzare} i task critici rispetto a quelli
best-effort, e quale sia l'\emph{overhead} del monitoraggio sui tempi di risposta. La valutazione
si basa su una campagna di 485 esecuzioni misurate su una piattaforma resa deterministica in modo
controllato. Il codice strumentato \`e stato fornito come base di partenza; il nostro contributo
riguarda la costruzione della piattaforma di misura, il disegno degli esperimenti, la loro
esecuzione e l'analisi dei risultati.
\end{minipage}
\end{center}

\vspace{1em}
{\small\tableofcontents}
\newpage

\section{La piattaforma sperimentale}

Misurare l'overhead di una libreria di tracing significa cercare differenze dell'ordine delle
decine di microsecondi. Su una macchina non preparata, la varianza dovuta a frequenza variabile,
migrazioni di processo e interruzioni \`e di gran lunga superiore al segnale che si vuole
osservare. La prima parte del lavoro \`e stata quindi dedicata a rendere la piattaforma
sufficientemente stabile da rendere significative le misure successive.

\subsection{Hardware e sistema}

Le misure sono state eseguite su un portatile con processore AMD Ryzen 7 3700U (4 core fisici,
8 thread SMT, TDP 15\,W) e kernel Linux \texttt{6.12.79-rt17} con patch PREEMPT-RT, installato
in dual boot su hardware reale e non in virtualizzazione.

Il processore espone 8 CPU logiche appaiate a due a due sui 4 core fisici:

\begin{center}\ttfamily\small
cpu0,\,cpu1 $\rightarrow$ core 0 \quad
cpu2,\,cpu3 $\rightarrow$ core 1 \quad
cpu4,\,cpu5 $\rightarrow$ core 2 \quad
cpu6,\,cpu7 $\rightarrow$ core 3
\end{center}

Tutti i core condividono la stessa cache L3, per cui la scelta di \emph{quali} core utilizzare \`e
indifferente; conta invece che si tratti di core \textbf{interi}. Questo ha determinato due
decisioni.

\textbf{Il task critico e i task di disturbo devono stare su core fisici diversi.} Due thread SMT
dello stesso core condividono unit\`a di esecuzione, cache L1 e L2: collocandoli l\`i,
l'interferenza misurata sarebbe stata in buona parte contesa hardware, cio\`e esattamente il
fenomeno che lo studio deve tenere separato dagli effetti di scheduling e di telemetria. Il task
critico \`e stato quindi assegnato a \texttt{cpu2} (core 1) e i task di disturbo a \texttt{cpu6}
(core 3).

\textbf{Non si isola una CPU lasciando fuori il suo gemello SMT.} Isolando ad esempio le sole
\texttt{cpu2} e \texttt{cpu4}, il carico di sistema che continua a girare sui rispettivi gemelli
sottrarrebbe risorse alle CPU nominalmente isolate, che risulterebbero tali solo sulla carta.
L'isolamento comprende quindi i due core interi: \texttt{2,3,6,7}.

\subsection{Parametri di avvio del kernel}

La riga di comando del kernel \`e stata configurata come segue:

\begin{lstlisting}
isolcpus=managed_irq,domain,2,3,6,7  nohz_full=2,3,6,7
rcu_nocbs=2,3,6,7  irqaffinity=0,1,4,5
\end{lstlisting}

Ciascun parametro risponde a una sorgente di disturbo distinta.

\begin{itemize}
\item \texttt{isolcpus=domain} rimuove le CPU indicate dai domini di bilanciamento dello
scheduler: nessun processo vi viene migrato automaticamente, e vi finisce solo ci\`o che viene
collocato esplicitamente.

\item \texttt{isolcpus=managed\_irq} agisce sul lato di sottomissione delle richieste di I/O: il
livello a code multiple del sottosistema a blocchi non utilizza le code il cui interrupt \`e
associato a una CPU isolata. \`E utile chiarire un equivoco: questo parametro \textbf{non} modifica
l'affinit\`a degli interrupt gestiti dal kernel, che rimane invariata; l'effetto \`e che quelle
code non vengono pi\`u utilizzate, e la verifica corretta \`e quindi il conteggio degli interrupt,
non l'affinit\`a dichiarata. Sulla nostra macchina le code del disco che ricadrebbero sulle CPU
isolate mostrano contatori a zero.

\item \texttt{nohz\_full} elimina il tick periodico di scheduling sulle CPU indicate quando vi
gira un solo task eseguibile. L'effetto \`e stato misurato e non assunto: in 5 secondi di
esecuzione all'interno dell'area isolata si contano \textbf{106} interruzioni di temporizzazione
locale sulla CPU del task critico, contro le 5000 che si avrebbero con il tick pieno alla
frequenza configurata di 1000\,Hz, e contro le 623--1216 osservate sulle CPU di servizio.

\item \texttt{rcu\_nocbs} sposta l'esecuzione dei callback di RCU fuori dalle CPU isolate.

\item \texttt{irqaffinity} fa s\`i che gli interrupt \emph{nascano} affini alle sole CPU di
servizio. Senza questo parametro, alcune periferiche (controller grafico, USB, audio, tastiera)
allocano i propri interrupt sulle CPU isolate, e occorre riassegnarli a posteriori. Dopo la
verifica, tali interrupt risultano affini a \texttt{0-1,4-5} e presentano conteggio nullo sulle
CPU isolate gi\`a prima di attivare l'isolamento a runtime.
\end{itemize}

Restano due categorie non eliminabili, di cui si \`e tenuto conto: i thread di kernel per-CPU
(migrazione, softirq, timer, RCU), che esistono per costruzione su ogni CPU e di cui l'isolamento
riduce l'attivit\`a ma non l'esistenza; e un piccolo numero di interrupt gestiti direttamente dal
kernel, la cui affinit\`a non \`e modificabile ma che, come verificato, non vengono sollecitati.

\subsection{Isolamento a runtime}

Ai parametri di avvio si affianca un isolamento applicato a ogni sessione di misura, che crea due
insiemi di CPU distinti: uno di servizio (\texttt{0,1,4,5}) in cui vengono confinati tutti i
processi esistenti, e uno riservato (\texttt{2,3,6,7}) in cui gira soltanto ci\`o che viene
lanciato esplicitamente. Lo stesso passo riassegna alle CPU di servizio le affinit\`a degli
interrupt che un driver avesse eventualmente allocato altrove dopo l'avvio.

L'efficacia \`e stata verificata sperimentalmente confrontando lo stesso carico dentro e fuori
l'area isolata. Il jitter del periodo, definito come differenza fra periodo massimo e mediano,
passa da 1420--2418\us{} all'esterno a 51--180\us{} all'interno, con un miglioramento di circa un
ordine di grandezza. Sottoponendo la macchina a carico artificiale sulle CPU di servizio, i valori
misurati all'interno dell'area isolata restano indistinguibili da quelli a macchina scarica,
mentre all'esterno il periodo mediano cresce del 40\,\% con un terzo delle iterazioni perse.

\subsection{Stabilit\`a della frequenza}

Il kernel real-time in uso \`e compilato senza il sottosistema di gestione della frequenza e senza
quello degli stati di riposo del processore. Non esistono quindi i consueti governor, e gli
strumenti abituali non sono applicabili: le transizioni di frequenza sono decise dal firmware e il
sistema operativo pu\`o soltanto osservarle. Il rovescio positivo \`e che, non essendovi stati di
riposo gestiti dal kernel, non si osservano picchi di latenza al risveglio; inoltre il contatore
temporale del processore \`e indipendente dalla frequenza, per cui la misura del tempo resta
corretta.

L'unica leva disponibile agisce sui registri di controllo del processore, ed \`e stata utilizzata
per disabilitare il meccanismo di sovra-frequenza automatica e per richiedere lo stato di
prestazione massimo. L'effetto \`e stato quantificato:

\begin{table}[h]\centering\small
\begin{tabular}{lcc}
\toprule
 & \textbf{sovra-frequenza attiva} & \textbf{frequenza fissata} \\
\midrule
frequenza osservata sui core & 2032--2670\,MHz (dispersione 31\,\%) & 2229--2296\,MHz (2.9\,\%) \\
costo per iterazione stimato & 18--21\,ns (dispersione 17\,\%) & 29--30\,ns (3.4\,\%) \\
durata della fase di lavoro & $+29$/$+44$\,\% sul valore richiesto & $+6$/$+9$\,\% \\
periodo, 99\textsuperscript{o} percentile & 11572--12305\us & 11165--11379\us \\
\bottomrule
\end{tabular}
\caption{Effetto della fissazione della frequenza, misurato su esecuzioni ripetute.}
\end{table}

La dispersione del 31\,\% sulla frequenza si trasferisce direttamente sulle misure temporali, e
sarebbe stata di gran lunga superiore all'overhead da misurare. Con la frequenza fissata, il valore
osservato su tutte le 485 esecuzioni della campagna \`e risultato costante a \textbf{2295\,MHz},
con temperatura compresa fra 48 e 57\,\textdegree C e nessun intervento di limitazione termica.

Le scritture su tali registri non sopravvivono al riavvio: vanno ripetute a ogni sessione, e il
loro esito viene verificato automaticamente prima di ogni campagna, che si interrompe se la
condizione non \`e soddisfatta.

\subsection{Costante di calibrazione}

Il generatore di carico determina la quantit\`a di lavoro da eseguire dividendo la durata
richiesta per un costo unitario stimato. Tale costo pu\`o essere misurato all'avvio oppure fissato
nella configurazione. Abbiamo scelto di fissarlo, per due ragioni.

La prima \`e di determinismo: la misura automatica esegue attese ripetute fino a convergenza e
richiede alcuni secondi non deterministici a ogni esecuzione, ma soprattutto restituisce un valore
che dipende dallo stato istantaneo della macchina. Un valore diverso dell'1\,\% sposta la quantit\`a
di lavoro effettivamente eseguito del 3.5\,\%, cio\`e molto pi\`u dell'effetto che la campagna deve
misurare.

La seconda \`e di comparabilit\`a: con la costante fissata, tutte le celle dell'esperimento
eseguono esattamente lo stesso numero di iterazioni del ciclo di carico, e le differenze osservate
sono attribuibili al solo fattore in studio. Due campagne indipendenti hanno stimato il costo reale
in circa 28.8\,ns per iterazione; poich\'e il parametro accetta solo valori interi, la scelta
ricade fra 28 (che darebbe un errore del $+2.7$\,\% sulla durata nominale) e 29 ($-0.8$\,\%). \`E
stato adottato il valore \textbf{29} su tutte le esecuzioni.
